\documentclass[11pt]{article}
\usepackage[T1]{fontenc}
\usepackage{lmodern,amsmath,amssymb,amsthm}
\usepackage[margin=1in]{geometry}
\usepackage[hidelinks]{hyperref}
\newtheorem{theorem}{Theorem}
\newcommand{\PP}{\mathbb P}
\newcommand{\OO}{\mathcal O}
\newcommand{\GK}{\operatorname{GKdim}}
\title{A projective-line quotient of the Rogalski--Sierra algebra}
\author{Alec Kriebel\\
\small Independent researcher\\
\small \href{https://orcid.org/0009-0001-9320-500X}{ORCID: 0009-0001-9320-500X}}
\date{16 October 2026 \quad Version 1.0}
\begin{document}
\maketitle
\begin{abstract}
For the six-relation presentation \(R(\alpha,\beta)\) in Rogalski and
Sierra's 2010 report and 2012 article, the linear ideal
\((x_3-x_1,x_4-x_2)\) has a projective-line twisted homogeneous coordinate
ring as its exact quotient. It is a nonzero homogeneous completely prime
ideal strictly below the augmentation ideal. This gives negative answers
to both printed assertions: uniqueness of the nontrivial graded prime
and absence of nontrivial maps to twisted homogeneous coordinate rings.
The calculation works over every field whenever \(\beta^2\ne1\), hence
for all the prescribed algebraically independent parameters.
The classical twist machinery is credited; historical priority of this
explicit quotient remains unresolved.
\end{abstract}

\section{The displayed problem}
Let \(k\) be a field, let \(\alpha,\beta\in k\), and give the generators
degree one in
\[
 R=R(\alpha,\beta)=k\langle x_1,x_2,x_3,x_4\rangle/(f_1,\ldots,f_6),
\]
where
\begin{align*}
 f_1&=x_1(\alpha x_1-x_3)+x_3(x_1-\alpha x_3),\\
 f_2&=x_1(\alpha x_2-x_4)+x_3(x_2-\alpha x_4),\\
 f_3&=x_2(\alpha x_1-x_3)+x_4(x_1-\alpha x_3),\\
 f_4&=x_2(\alpha x_2-x_4)+x_4(x_2-\alpha x_4),\\
 f_5&=x_1(\beta x_1-x_2)+x_4(x_1-\beta x_2),\\
 f_6&=x_1(\beta x_3-x_4)+x_4(x_3-\beta x_4).
\end{align*}
In particular, the left factor in the second summand of both \(f_5\)
and \(f_6\) is \(x_4\).

The official report~\cite[pp.~1361--1362]{OWR} gives these relations and
asks whether, for algebraically independent \(\alpha,\beta\) over the
prime subfield of an uncountable algebraically closed \(k\), the
augmentation ideal \(R_+=\bigoplus_{n>0}R_n\) is the only nontrivial
graded prime ideal, and whether \(R\) has no nontrivial homomorphism to
any twisted homogeneous coordinate ring. Both assertions recur in the
published article~\cite[p.~1198]{RS}, whose Theorem~3.5(2) on p.~1202
gives this presentation with parameters \(c,d\) in place of
\(\alpha,\beta\). The printed assertions impose no injectivity,
birationality or lower bound on the dimension of the image.
We exhibit a unital graded surjection with infinite-dimensional image
of Gelfand--Kirillov dimension two.

\section{The quotient}
We use positive pullbacks in the convention
\[
 B(X,L,\sigma)_n=H^0(X,L_n),\qquad
 L_n=L\otimes\sigma^*L\otimes\cdots\otimes(\sigma^{n-1})^*L,
 \qquad f\star g=f\otimes(\sigma^m)^*g
\]
for \(f\) of degree \(m\), with \(L_0=\OO_X\). This is the classical
twisted homogeneous coordinate ring construction~\cite{AV, RZ}.

\begin{theorem}
For every field \(k\), every \(\alpha\in k\), and every \(\beta\in k\)
with \(\beta^2\ne1\), put
\[
 I=(x_3-x_1,x_4-x_2)\subset R,\qquad
 \sigma[s:t]=[s+\beta t:\beta s+t].
\]
Then
\[
 R/I\ \cong\
 k\langle X,Y\rangle/(\beta X^2-XY+YX-\beta Y^2)
 \ \cong\ B(\PP^1_k,\OO(1),\sigma).
\]
The isomorphism sends \(x_1,x_3\) to \(s\) and \(x_2,x_4\) to \(t\).
Moreover, \(0\ne I\subsetneq R_+\), \(I\) is homogeneous and completely
prime, and \(R/I\) has Hilbert series \((1-z)^{-2}\) and
\(\GK(R/I)=2\).
\end{theorem}

\begin{proof}
Imposing \(x_3=x_1=X\) and \(x_4=x_2=Y\) sends the six relations to
\(0,0,0,0,Q,Q\), where
\[
 Q=\beta X^2-XY+YX-\beta Y^2.
\]
The universal property of free-algebra quotients therefore gives the
first displayed isomorphism, with exactly this defining ideal.

Let \(S=k[s,t]\) and let \(\theta\) be the graded algebra automorphism
\[
 \theta(s)=s+\beta t,\qquad \theta(t)=\beta s+t.
\]
It is invertible since its determinant is \(1-\beta^2\).
On the graded vector space \(S\), define
\(f\star g=f\theta^m(g)\) for homogeneous \(f\in S_m\), and extend
bilinearly. The automorphism law proves associativity and the unit
property. A product of two nonzero homogeneous elements is nonzero;
taking highest nonzero homogeneous components proves that this
algebra \(B\) is a domain for arbitrary elements as well.
Its degree \(n\) piece has dimension \(n+1\).
It is generated by degree one: inductively,
\(S_n\theta^n(S_1)=S_{n+1}\).

The linear lift of \(\sigma\) above identifies \(\sigma^*\) on
\(H^0(\OO(1))\) with \(\theta\), and identifies \(L_n\) with \(\OO(n)\).
Under these compatible identifications the section multiplication is
exactly \(\star\). Thus \(B=B(\PP^1,\OO(1),\sigma)\).
The line bundle is \(\sigma\)-ample: \(L_n\cong\OO(n)\), so Serre
vanishing applies to every coherent sheaf tensored with \(L_n\).

Direct computation in the commutative polynomial ring \(S\) gives
\[
 \beta\,s\star s-s\star t+t\star s-\beta\,t\star t=0.
\]
Since \(s,t\) generate \(B\), there is a graded surjection
\(A:=k\langle X,Y\rangle/(Q)\twoheadrightarrow B\).
To prove injectivity without any confluence assumption, use
\(U=X-Y,\ V=Y\). This change of generators is invertible in every
characteristic, and
\[
 Q=\beta U^2+(\beta-1)UV+(\beta+1)VU.
\]
Because \(\beta+1\ne0\), the relation rewrites \(VU\) as
\[
 VU=\frac{1-\beta}{1+\beta}UV-\frac{\beta}{1+\beta}U^2.
\]
Every replacement strictly decreases fixed-length lexicographic order
with \(V>U\). Reduction therefore terminates and shows that
\(U^iV^{n-i}\), \(0\le i\le n\), span \(A_n\).
Consequently \(\dim A_n\le n+1=\dim B_n\). The surjection forces equality
in each degree, hence is an isomorphism and the kernel of \(R\to B\)
is exactly \(I\).

The relations of \(R\) are quadratic, so its four degree-one generators
are linearly independent. In particular \(x_3-x_1\ne0\).
The ideal \(I\) is homogeneous, two-sided and contained in \(R_+\).
Since \(x_1\) maps to the nonzero element \(s\), it is not in \(I\);
thus \(0\ne I\subsetneq R_+\). The domain quotient proves complete
primality, and hence both ordinary and graded primality.
Finally \(\dim B_n=n+1\), so the Hilbert series is \((1-z)^{-2}\);
the cumulative dimension is \((n+1)(n+2)/2\), proving \(\GK B=2\).
\end{proof}

Algebraic independence of \(\alpha,\beta\) implies \(\beta^2\ne1\).
Thus the ideal and surjection refute both assertions in the stated
parameter regime. The image comes from a positive-dimensional
projective curve; it is neither an augmentation to \(k\) nor a
finite-dimensional image.

\section{Scope and historical context}
No classification of all primes is claimed. Nor does a quotient
contradict the separate obstruction to sharing a graded quotient ring
with a twisted homogeneous coordinate ring discussed in~\cite[p.~1198]{RS}.
The source's surface algebra and this curve quotient have different
Gelfand--Kirillov dimensions.

The theorem concerns the displayed presentation in all characteristics.
The journal's geometric identification uses
\(c=(\vartheta-1)/(\vartheta+1)\) and
\(d=(\rho-1)/(\rho+1)\)~\cite[Theorem~3.5]{RS}.
These preserve independent parameters in characteristic different from
two; in characteristic two the fractions collapse to one when defined.
The theorem does not identify the independent-parameter
characteristic-two presentation with that geometric family or transfer
the latter's noetherian or dimension results.
At the excluded values \(\beta=\pm1\), the corresponding \(Q\) factors
as \((X+Y)(X-Y)\) or \(-(X-Y)(X+Y)\); the quotient then has nonzero
linear zero divisors. At \(\beta=0\) it is the commutative polynomial
plane, consistent with the stronger theorem.

The two-generator algebras, degree-reduction argument and polynomial
graded twists are classical. The source-specific point is the
substitution \((x_1,x_2,x_3,x_4)\mapsto(X,Y,X,Y)\) and its exact
projective-line quotient.
Later Nevins--Sierra preprints~\cite[v1, p.~25; v2, pp.~28--29]{NS}
repeat the contrary no-map assertion for the Rogalski--Sierra algebra.
Their Theorem~1.2 requires additional point-tail moduli and geometric
hypotheses, including dimension at least two; a universal factor in
Theorem~4.10 need not itself be a twisted homogeneous coordinate ring.
Likewise the canonical-map theorem in~\cite{RZ} requires strong
noetherianity, which the generic source family lacks
\cite[p.~1236]{RS}. These theorems alone do not supply an earlier
resolution. Campbell's thesis~\cite[Chapter~6 and \S8.2]{Campbell}
treats infinitesimal deformations at a baseline algebra and function
skew fields of another deformation family; those statements are not
this exact generic algebra's homomorphic curve quotient.

Two independently positioned priority audits verified no full matching
predecessor in the inspected primary passages. Historical priority
remains unresolved. Unread journal versions, older papers and books,
unindexed work and incomplete citation coverage remain gaps.
In particular, the repeated assertion above is attributed to the
inspected arXiv versions, not to an unread journal body.
No firstness, exhaustive-absence or current worldwide-open claim
is made.

\section*{Verification, disclosure and availability}
The companion source archive preserves all 58 immutable candidate files
and three original archives, exact selected-record and statement
identities, two independent mathematical approaches and two independent
priority audits. It includes standard-library Python verification and
finite diagnostics, including characteristic-two and singular-parameter
controls. The mathematical audits use both a terminating spanning
argument and an independent tensor-transform description of the
commutator ideal. These computations support the written universal
proof; they are not sampling proofs or proof-assistant certification.
Full third-party PDFs, extracts and page images are not redistributed.
The audit materials and subsequent whole-package reviews are recorded
at \href{https://github.com/AlecKriebel/Math/pull/789}{repository PR~789}.

AI tools were used extensively in mathematical investigation,
literature search, drafting, computation, reproduction and adversarial
review. This is an unrefereed preprint, without external human peer
review or formal proof-assistant verification.
The author accepts responsibility for the claims.
The note and authored companion materials are licensed
\href{https://creativecommons.org/licenses/by/4.0/}{CC BY 4.0}.

\begin{thebibliography}{9}
\bibitem{OWR}
S.~J.~Sierra, \emph{Birationally commutative projective surfaces of
GK-dimension 4} (joint work with D.~Rogalski), in
\emph{Oberwolfach Report 22/2010},
pp.~1360--1362.
\href{https://doi.org/10.4171/OWR/2010/22}{doi:10.4171/OWR/2010/22}.
\bibitem{RS}
D.~Rogalski and S.~J.~Sierra,
\emph{Some projective surfaces of GK-dimension 4},
Compositio Math.~148 (2012), 1195--1237.
\href{https://doi.org/10.1112/S0010437X12000188}{doi:10.1112/S0010437X12000188}.
\bibitem{AV}
M.~Artin and M.~Van den Bergh,
\emph{Twisted homogeneous coordinate rings},
J.~Algebra 133 (1990), 249--271.
\bibitem{RZ}
D.~Rogalski and J.~J.~Zhang,
\emph{Canonical maps to twisted rings},
Math.~Z.~259 (2008), 433--455;
\href{https://arxiv.org/abs/math/0409405v2}{arXiv:math/0409405v2}.
\bibitem{NS}
T.~A.~Nevins and S.~J.~Sierra,
\emph{Na\"ive blowups and canonical birationally commutative factors},
\href{https://arxiv.org/abs/1206.0760v1}{arXiv:1206.0760v1} (2012),
\href{https://arxiv.org/abs/1206.0760v2}{v2} (arXiv revision 2014;
retrieved PDF internally dated 15 November 2018).
\bibitem{Campbell}
C.~J.~M.~Campbell,
\emph{Deformation theory of a birationally commutative surface of
Gelfand--Kirillov dimension 4}, Ph.D. thesis, University of Edinburgh,
2016. \href{https://hdl.handle.net/1842/22886}{hdl:1842/22886}.
\end{thebibliography}
\end{document}
